%==============================================================================
% Capitulo 6: Implementacion
%==============================================================================
\section{Implementacion}

\subsection{Entorno de despliegue y dependencias}

La implementación del sistema se llevo a cabo utilizando exclusivamente tecnologías de código abierto con el objetivo de garantizar la portabilidad, la reproducibilidad y la auditoria de seguridad de la solución en organizaciones medianas. Toda la infraestructura se encapsula mediante contenedores Docker, orquestados localmente por Docker Compose para los entornos de desarrolló y preproduccion, y preparados para migración a un cluster Kubernetes~1.30 para el entorno productivo.

Las versiones específicas de los componentes son: N8N~1.62, Python~3.12, PostgreSQL~15.5 con volumen persistente, FastAPI~0.115, Uvicorn~0.32, SQLAlchemy~2.0, y el cliente google-generativeai~0.8 para la comunicación con Gemini~2.5 Flash. Todos los servicios se despliegan sobre imagenes Docker oficiales (\texttt{python:3.12-slim} y \texttt{postgres:15.5-alpine}). La gestión de dependencias del modulo Python se realiza mediante un archivo \texttt{requirements.txt} con versiones fijadas (\emph{pinned}) para garantizar la reproducibilidad de las compilaciones en distintos entornos.

\subsection{Construccion del modulo Python}

La selección del marco de trabajo para el servicio de procesamiento implico la evaluación de tres alternativas: Flask, Django REST Framework y FastAPI. Flask, ampliamente adoptado en el ecosistema Python por su simplicidad, carece de validación nativa de esquemas y de generación automática de documentación OpenAPI, funcionalidades que el proyecto requeria para garantizar la calidad de la interfaz REST. Django REST Framework, por su parte, ofrece un conjunto completó de herramientas para la construcción de interfaces de programacion de aplicaciones, pero su modelo de operación sincronico y la inclusion de componentes no requeridos por el proyecto ---tales como el ORM de Django y el motor de plantillas--- incrementaban la complejidad de la base de código sin aportar beneficios funcionales para una arquitectura de microservicio acotado. FastAPI~0.115, sustentado sobre el servidor ASGI Uvicorn~0.32, ofrece tres ventajas determinantes para el proyecto: validación automática de esquemas de entrada y salida mediante Pydantic~v2, generación de documentación interactiva conforme a la especificación OpenAPI~3.1 sin configuración adicional, y operación asincronica nativa que permite manejar las llamadas concurrentes desde N8N sin bloquear el hilo de ejecución principal \parencite{Tiangolo2024_fastapi}. La madurez del ecosistema de FastAPI, su adopción creciente en la industria y la disponibilidad de documentación exhaustiva en espanol consolidaron la decision.

El modulo Python expone su contrato REST organizado en capas, según la estructura descrita en el Capitulo~5, con una separación clara entre responsabilidades:

\begin{itemize}
  \item \textbf{models}: entidades del modelo de dominio mapeadas mediante SQLAlchemy~2.0.
  \item \textbf{schemas}: modelos de validación y serializacion con Pydantic~v2.
  \item \textbf{services}: lógica de aplicación (IncidenteService, ClasificacionService).
  \item \textbf{repositories}: acceso a datos, aislando las consultas SQL del resto de la aplicación.
  \item \textbf{classifiers}: lógica del clasificador hibrido (filtro deterministico + LLM).
  \item \textbf{routes}: enrutadores REST que delegan en los servicios correspondientes.
  \item \textbf{core}: configuración de base de datos, manejo de errores y logging estructurado.
\end{itemize}

La construcción se realiza bajo un esquema de inyeccion de dependencias gestionado por las funcionalidades nativas de FastAPI (\texttt{Depends}), lo que facilita el reemplazo de implementaciones durante la ejecución de pruebas unitarias y de integración. Esta arquitectura permite, por ejemplo, sustituir el clasificador real por una implementación simulada durante las pruebas sin modificar el código de los servicios que lo consumen. La gestión de credenciales se realiza exclusivamente mediante variables de entorno, en línea con el principio de configuración de la metodología Twelve-Factor App \parencite{Wiggins2017_12factor}.

\subsection{Construccion del flujo en N8N}

El flujo de orquestación implementado en N8N comprende 12 nodos principales encadenados secuencialmente. La eleccion de N8N como plataforma de orquestación, frente a alternativas como Apache Airflow, Zapier o Make, obedecio a tres criterios concurrentes. El primero, de naturaleza operativa, radica en que N8N admite despliegue completamente autoalojado bajo Docker, manteniendo los flujos de procesamiento y los datos dentro del perimetro organizacional. Zapier y Make operan exclusivamente como servicios en la nube, lo que las descarta en escenarios donde el procesamiento de datos sensibles exige control de infraestructura. Apache Airflow, aunque autoalojable, disena sus flujos como grafos aciclicos dirigidos programados por intervalos temporales; ese paradigma resulta adecuado para procesos de extracción, transformación y carga de datos pero forzadamente inadecuado para un sistema reactivo que debe responder en tiempo real a eventos de entrada no programados \parencite{Apache2024_airflow}. El segundo criterio refiere al modelo de construcción visual que N8N ofrece sobre un editor de nodos arrastrables: la curva de aprendizaje de esta interfaz es sustancialmente menor que la de alternativas centradas en definición programatica de flujos, lo que permitió al equipo concentrar el esfuerzo en la lógica de negocio antes que en la sintaxis de definición de procesos. El tercer criterio se vincula con la licencia de código fuente disponible (\emph{fair-code}) que habilita la auditoria de seguridad del nucleo de orquestación \parencite{n8n2024_docs}.

La composición de los 12 nodos responde a un patron de procesamiento en cinco etapas que se replica para cada canal:

\begin{enumerate}[leftmargin=*]
  \item \textbf{Recepcion}: Tres disparadores (\emph{triggers}) reciben las entradas desde correo electrónico (IMAP), formulario web (Webhook) y telefonia (Webhook post-transcripción). La decision de separar los disparadores por canal, en lugar de unificar la recepción mediante un único webhook con discriminacion interna, se fundamenta en la necesidad de preservar la configuración de autenticacion y los parámetros de reintento específicos de cada protocolo.
  \item \textbf{Normalizacion}: Un nodo de transformación (\emph{Function}) homogeniza la estructura del mensaje en un formato unificado con campos de identificador único (UUID v4), marca temporal, canal de origen, remitente y descripción textual normalizada. Esta etapa elimina las diferencias de formato entre canales y garantiza que las capas posteriores procesen siempre la misma estructura de datos, independientemente de la via de entrada.
  \item \textbf{Clasificacion}: Un nodo HTTP invoca el endpoint \texttt{POST /api/v1/clasificar} del modulo Python enviando la descripción textual y recibiendo la categoría sugerida y la confianza asociada en formato JSON.
  \item \textbf{Evaluacion}: Un nodo condicional (\emph{IF}) evalua la confianza devuelta por el clasificador; si esta por debajo del umbral de 0,70, deriva el incidente a la cola de revisión humana y notifica al operador designado. Si la confianza supera el umbral, el flujo continua automáticamente.
  \item \textbf{Persistencia y notificacion}: Un nodo HTTP invoca \texttt{POST /api/v1/incidentes} para persistir el ticket en la base de datos, y tres nodos paralelos notifican al usuario por el canal correspondiente y registran la ejecución en el sistema de auditoria. El paralelismo de la notificacion garantiza que la demora de un canal (por ejemplo, el envio de correo electrónico) no bloquee el registro de auditoria ni la confirmacion por otros medios.
\end{enumerate}

La trazabilidad completa de cada ejecución se preserva almacenando los identificadores de nodo, las marcas temporales de entrada y salida de cada etapa, y los datos de entrada y salida de cada paso del flujo, generando un historial auditable que resulta especialmente valioso en contextos donde la Ley~25.326 exige demostrar el camino recorrido por los datos personales \parencite{Ley25326_2000}. La exportación completa del flujo en formato JSON se incluye en el Anexo~E.

\subsection{Integracion del canal telefonico}

La integración con Twilio se realizó mediante el producto Programmable Voice, complementado con la transcripción automática nativa. El flujo telefonico opera de la siguiente manera: el usuario marca un número virtual asignado a la organización; un script TwiML reproduce un mensaje de bienvenida en espanol rioplatense y solicita al usuario describir su problema durante un máximo de 45 segundos; la grabacion finaliza cuando el usuario presiona la tecla numeral (\texttt{\#}) o transcurrido el tiempo máximo; la grabacion se envia al servicio de transcripción de Twilio; y, una vez transcrito el contenido, Twilio invoca un webhook expuesto por N8N con la transcripción textual. A partir de ese punto, el flujo continua de manera identica a la de los canales escritos.

La latencia total del canal telefonico ---medida desde el cuelgue del usuario hasta la confirmacion del número de incidente--- se ubico en un rango medio de 12 a 15 segundos según las observaciones operativas durante la fase de validación. Esta latencia incluye: transcripción (~5--7 s), clasificación (~2--3 s para incidentes resueltos por el filtro deterministico, ~5--8 s cuando se consulta al LLM) y persistencia (~1 s).

\subsection{Desafios de integracion}

La implementación del sistema enfrento tres desafios de integración que merecen documentación explicita por su potencial utilidad en replicaciones futuras.

El primer desafio se presentó en la comunicación asincronica entre N8N y el servicio de clasificación. El nodo HTTP de N8N invoca el endpoint de clasificación de manera sincronica y espera la respuesta del servidor antes de continuar el flujo. Esta arquitectura, aunque simple de implementar, introduce un punto de fallo único: si el servicio Python no responde dentro del tiempo de espera configurado en el nodo HTTP, el flujo completó se detiene y el incidente queda sin procesar. La solución consistio en implementar un mecanismo de reintento con retroceso exponencial en el propio nodo HTTP de N8N, combinado con una cola de respaldo en el servicio Python que retiene los mensajes recibidos durante ventanas de saturacion y los procesa en orden cuando la carga se normaliza. Esta combinación de reintento en el cliente y cola en el servidor elimino las perdidas de incidentes observadas durante las pruebas de carga iniciales.

El segundo desafio surgio de la heterogeneidad de los formatos de entrada. Los correos electrónicos incluyen encabezados, firmas corporativas y reenvios que contaminan el texto que el clasificador recibe; las llamadas telefonicas transcritas introducen errores de reconocimiento del habla, particularmente en terminos técnicos (por ejemplo, \texttt{API} transcrito como \texttt{a p i} o \texttt{apeye}); y los formularios web, aunque proveen texto limpio, llegan con una proporción significativa de campos vacios o con descripciones de una sola palabra. La respuesta a esta heterogeneidad fue doble. Por un lado, se implementó un preprocesador en el nodo de normalización de N8N que elimina firmas, encabezados de reenvio y líneas de disclaimer mediante expresiones regulares específicas del dominio de correo institucional. Por el otro, se enriquecio el prompt enviado al clasificador con instrucciones explicitas para interpretar terminos técnicos frecuentemente mal transcritos, lo que elevo la robustez de la clasificación sobre transcripciones automáticas en aproximadamente 4 puntos porcentuales de F1 respecto de la configuración inicial sin estas instrucciones.

El tercer desafio, de naturaleza operativa, se vincula con la persistencia transaccional de la sesion de clasificación. El sistema debe garantizar que una vez que el clasificador asigna una categoría, el incidente se persiste con esa decision o, en caso de falla de la base de datos, el estado interno del sistema refleje la inconsistencia y pueda recuperarse. La solución adopto un enfoque pragmatico: N8N ejecuta la clasificación y la persistencia como pasos secuenciales dentro del mismo flujo; si la persistencia falla, el flujo se reintenta desde el paso de clasificación, lo que garantiza que el incidente eventualmente se registre aun a costa de duplicar la llamada al clasificador. La tabla de trazabilidad (\texttt{clasificacion\_log}) fue diseñada para admitir multiples registros de clasificación por incidente, preservando así el historial de cada intento.

\subsection{Conexion con el proceso Scrumban}

La construcción del sistema fue organizada bajo el esquema Scrumban descrito en el Capitulo~4, ejecutando seis iteraciones de dos semanas cada una. La vinculacion entre cada iteracion y los artefactos de implementación producidos se detalla a continuación:

\begin{enumerate}[label=Iteracion \arabic*:, leftmargin=*]
  \item La elicitacion de requisitos (Semanas 1--2) produjo el diseñó preliminar de la arquitectura de cinco capas y la especificación de la interfaz REST, documentadas en el Capitulo~5. El entregable concreto fue el documento de arquitectura y la definición del contrato OpenAPI inicial.
  \item La construcción del modulo Python (Semanas 3--4) entrego el modelo de datos con sus cinco entidades, los repositorios con operaciones CRUD asincronicas y los servicios de negocio. Al cierre de esta iteracion, el modulo Python exponia el endpoint de salud y aceptaba la creacion de incidentes mediante \texttt{POST /api/v1/incidentes}, con cobertura de pruebas unitarias superior al 80~\%.
  \item La configuración inicial del flujo N8N (Semanas 5--6) implementó el canal de correo electrónico como via prioritaria, estableciendo el patron de cinco etapas que luego se replicaria para los canales restantes. Al finalizar esta iteracion, el sistema procesaba incidentes por correo electrónico de extremo a extremo, desde la recepción hasta la notificacion.
  \item La integración del canal telefonico y la persistencia (Semanas 7--8) incorporo Twilio Programmable Voice, el script TwiML y la configuración del webhook de transcripción. Esta iteracion fue la que mas expuso los desafios de heterogeneidad de formatos descritos en la Seccion~6.5, y su resolución consumio aproximadamente el 40~\% del esfuerzo de la iteracion.
  \item La validación funcional y el ajuste fino del clasificador (Semanas 9--10) se concentraron en la calibracion del filtro deterministico y del prompt de Gemini. Los ajustes se realizaron exclusivamente sobre incidentes representativos en entornos de preproduccion, sin exposición al corpus de validación de 200 casos, preservando la integridad de la etapa experimental posterior.
  \item La validación experimental y la documentación final (Semanas 11--12) ejecutaron el protocolo de pruebas descrito en el Capitulo~4, recolectaron las mediciones del corpus completó y produjeron el análisis estadístico presentado en el Capitulo~7.
\end{enumerate}

La trazabilidad entre iteraciones y artefactos, mediada por el tablero Kanban en GitHub Projects, permitió al equipo identificar tempranamente los desvios respecto del plan y redistribuir el esfuerzo cuando la complejidad de un componente superaba la estimacion inicial. La adopción de limites de trabajo en curso (WIP) por columna del tablero evito la acumulación de tareas a medio terminar, una práctica que resultó particularmente valiosa durante las iteraciones 4 y 5, donde la interdependencia entre modulos exigia finalizar componentes antes de integrarlos.

\subsection{Pruebas automatizadas}

La estrategia de pruebas se sustento en la piramide clasica descrita por \textcite{Crispin2009_agile_testing}. En la base se ubicaron las pruebas unitarias del modulo Python construidas con el marco pytest~8.3, alcanzando una cobertura del 87~\% del código de aplicación medida con coverage.py~7.6. Estas pruebas cubren los clasificadores (deterministico e hibrido), los servicios de negocio, los esquemas de validación y los repositorios de acceso a datos.

En el nivel intermedio se ejecutaron pruebas de integración que verificaron la interacción entre el modulo Python y la base de datos en una instancia de PostgreSQL desechable creada por contenedor en cada corrida de pruebas, utilizando SQLite en memoria para las pruebas unitarias y PostgreSQL real para las pruebas de integración. En el nivel superior se construyeron pruebas extremo a extremo que ejecutaban el flujo completó desde el envio de un correo simulado hasta la persistencia del ticket, validando la integración con N8N en su totalidad.

La integración continua (CI) se ejecuta automáticamente en cada solicitud de incorporacion (\emph{pull request}) al repositorio mediante GitHub Actions, ejecutando la suite completa de pruebas unitarias y de integración antes de permitir la fusion a la rama principal. El pipeline de CI incluye verificación de estilo de código (Ruff), ejecución de pruebas con reporte de cobertura, y verificación de sincronización de la especificación OpenAPI.

\newpage
